\documentclass[10pt,a4paper]{article}
\usepackage[a4paper,margin=8mm,top=8mm,bottom=8mm,head=0pt,foot=0pt,headheight=0pt,headsep=0pt,footskip=0pt]{geometry}
\usepackage{luatexja}
\usepackage{luatexja-fontspec}
\setmainjfont{Hiragino Sans}[BoldFont=Hiragino Sans W6]
\setsansjfont{Hiragino Sans}[BoldFont=Hiragino Sans W6]
\usepackage{longtable}
\usepackage{array}
\usepackage[table]{xcolor}

\pagestyle{empty}
\setlength{\parindent}{0pt}
\setlength{\parskip}{0pt}
\renewcommand{\baselinestretch}{0.90}

\definecolor{chead}{RGB}{25,80,160}
\definecolor{ckey}{RGB}{180,30,35}

\newcommand{\W}[1]{\textbf{#1}}
\newcommand{\K}[1]{\textcolor{ckey}{#1}}

\begin{document}
\footnotesize

% ===== HEADER NOTE =====
\noindent
\begin{minipage}{\textwidth}
\setlength{\fboxsep}{2pt}%
\fcolorbox{chead}{white}{%
\begin{minipage}{0.985\textwidth}
{\bfseries\color{chead}CCAR-P 用 一般英語ボキャブラリー（正答を左右する難単語 50）}\\[1pt]
\K{目的:} AI/IT 用語\K{以外}の、ビジネス・契約・アーキテクチャ寄りの硬い一般英語のなかで、問題文・選択肢・解説の意味を取り違えると\K{正答を落とす}語を精選。TOEIC 750+ のなかでも特に硬い（契約・法務・アーキ設計）帯を重視。\\
\K{出典:} Purcell 63問(purcell-pdf-CCAR-P.pdf)・Udemy D5/D6 演習(correct/incorrect/marked)より抽出。\\
\K{選択基準:} (1) AI/IT ジャーゴン除外 (2) 一般英語(正式・契約寄り) (3) TOEIC 750+ の硬め帯 (4) 正答ゲートとして作用する語。\\
\K{構成:} アルファベット順／見出し・品詞・意味・CCAR-P 文脈。文脈は独自の言い換え（出典文の引用なし）。
\end{minipage}}
\end{minipage}

\vspace{3pt}

% ===== VOCAB TABLE =====
\renewcommand{\arraystretch}{1.32}
\setlength{\tabcolsep}{3pt}

\begin{longtable}{@{}>{\raggedright\arraybackslash}p{23mm}
                    >{\raggedright\arraybackslash}p{13mm}
                    >{\raggedright\arraybackslash}p{38mm}
                    >{\raggedright\arraybackslash}p{114mm}@{}}
\rowcolor{chead!15}
\W{Word} & \W{POS} & \W{意味} & \W{CCAR-P 文脈（なぜ正答を左右するか）} \\
\endhead

\W{accountable} & adj. & 説明責任がある、責任を負う & 費用や品質の最終決定を誰が担うか。安価モデル強制をされた際、証拠を示した上で「責任ある決定者」に判断させるのが正解行動（上位への飛ばし越えは不正解）。 \\
\W{adversarial} & adj. & 敵対的な、攻撃的な & 規制領域の事前評価では安全境界やプロンプト注入を探る「敵対的テスト」が必須。これを含む選択肢が D4/D5 の正解。 \\
\W{artefact} & n. & 成果物、人工物 & ハンドオフ時の「成果物群」（意思決定記録・運用手順書・評価基線）が移管成功の鍵。英国綴り（artifact も可）。 \\
\W{breach} & n./v. & 違反、（契約・SLA の）侵害 & 広告基準や SLA の「侵害」。これを未然に防ぐか事後検出するかで、予防ガードレール・HITL・監視のいずれが正解かが決まる。 \\
\W{calibrate} & v. & 較正する、調整する & LLM 審査員を人間専門家の評価に定期的に「較正」する主観品質評価法。大規模ブランド評価での正解手法。 \\
\W{coarse} & adj. & 粗い、粒度が大きい & DAU 等の指標は「粗すぎる」ため品質劣化の早期検出に不適。leading indicator 選びで消去される選択肢。 \\
\W{consequential} & adj. & 重大な結果をもたらす & 重大な影響を持つ事項では人間への迂回路を示すべき。HITL 承認ゲートを置くべき位置の判断基準。 \\
\W{contractual} & adj. & 契約上の & SLA は「契約上の」応答時間要件。再ランク付け追加による遅延が SLA 内に十分余裕があるかで採用可否が決まる。 \\
\W{controller} & n. & （GDPR）統制者、個人情報取扱事業者 & データ取扱の指示を出す主体で処理者より上位の義務を負う。性能最適化が「誰の指示」に縛られるかの正解鍵。 \\
\W{defensible} & adj. & 擁護できる、正当化できる & モデル選定や成功基準が「説明可能（擁護できる）」であること。大規模トラフィックで最小モデルを選ぶ根拠。 \\
\W{defer} & v. & 先送りする、延期する & 決定や成功基準の「先送り」は概して不正解。データが揩っている場面で決断を遅らせる選択肢は避ける。 \\
\W{disparate} & adj. & 異質な、（差別的に）不均衡な & 保護群間の「不均衡な影響（disparate impact）」。集計精度が隠すバイアスを顕在化させる概念。 \\
\W{disaggregate} & v. & 細分化する、群別に分ける & 保護群別に評価指標を「群別化」すると集計値が隠すアルゴリズム・バイアスが顕在化する。公平性評価の正解手法。 \\
\W{disclaimer} & n. & 免責事項、注意書き & 毎日変わる情報へ「内容が異なる場合があります」と注記するだけでは不正解。鮮度問題の誤答肢の典型。 \\
\W{dogma} & n. & 教義、独断 & 「レイテンシ増は絶対に許されない」等は「独断」であり分析でない。SLA 対比での判断理由が不正解な理由。 \\
\W{escalate} & v. & （問題を）上位にエスカレートする & 対話前に上位委員会へ「エスカレート」するのは早計。証拠提示と合意形成が先の正解行動。 \\
\W{exemplar} & n. & 模範、実例 & 書式順守では完全な「実例（few-shot）」を示すのが最強の改善。記述による指示より具体的な目標として効く。 \\
\W{fabricate} & v. & 捏造する、でっち上げる & ツールが空振りした際に存在しない番号を「でっち上げる」のはハルシネーションの典型。原因診断の分類肢。 \\
\W{forfeit} & v. & 失い、没収される & 新モデル採用を半年遅らせると改善機会を「失う」。漸次展開ではなく待機を選ぶ肢の欠点として現れる。 \\
\W{incumbent} & adj./n. & 現行の、既存の & 自動化の費用を「現行の」手作業プロセスと比較する。費用比較の正解の枠組み。 \\
\W{irreversible} & adj. & 取り消せない、不可逆的な & HITL 承認ゲートは「不可逆的・高影響」な外部アクションの前に置く。設計判断の核心基準。 \\
\W{jurisdiction} & n. & 司法管轄権、法域 & 別「法域」の第三者分析プラットフォームへのデータ転送が GDPR 上の問題。越境データ取扱の鍵。 \\
\W{mandate} & n./v. & 命令、義務化する & コスト削減の「命令」に根拠なく従う、生成コードの全面手書きを「義務化」する等は不正解。 \\
\W{material} & adj. & 重大な、実質的な & 安価モデルの誤り率が「実質的に高い」、決定者に伝えるべき「重要な」情報。法的に意味のある差の基準。 \\
\W{mediate} & v. & 媒介する、仲介する & エージェント間連携では各エージェントが自組織のシステムへのアクセスを「仲介」する。他社ツールの直接公開を避ける設計。 \\
\W{minimisation} & n. & 最小化（データ最小化原則） & 目的に必要な個人情報だけを残す「データ最小化」。匿名化・削除で実装する GDPR の正解肢。英国綴り。 \\
\W{normative} & adj. & 規範的な & バイアスは測定可能な歪みだが、公平性は「規範的」な目標。「検出≠解決」を説明する概念。 \\
\W{obligation} & n. & 義務、（コンプライアンス）責務 & 健康データ等のコンプライアンス「義務」は設計段階で解決すべき。事後追加は違反であり改善でない。 \\
\W{premature} & adj. & 早すぎる、時期尚早の & バイアス検出＝解決済とする結論は「早計」。設計出力物を発見段階で決めるのも「時期尚早」。 \\
\W{prestige} & n. & 威信、名声 & 「最も高度なアーキテクチャ」は「名声」であって設計ドライバでない。ユースケース選定の不正解理由。 \\
\W{processor} & n. & （GDPR）処理者、受託者 & 統制者の指示に従う「処理者」。プロンプトキャッシュ等の最適化も指示と保存期限に縛られる。 \\
\W{proportionate} & adj. & 釣り合った、相当する & エラー処理はビジネスリスクに「見合った」ものに。過剰にも不足にもしない設計判断の基準。 \\
\W{proxy} & n. & 代理変数、代用指標 & 郵便番号等の「代理変数」が保護属性を暗に符号化しバイアスを生む。属性除去だけでは不十分な理由。 \\
\W{pseudonymise} & v. & 疑似匿名化する & データ境界外に出す前に個人データを「疑似匿名化」する。GDPR 対策の正解肢。英国綴り。 \\
\W{reconcile} & v. & 調和させる、両立させる & 性能最適化と統制者の保存期限・所在要件を「両立させる」判断が問われる。対立制約の処理。 \\
\W{redact} & v. & （機密部を）黒塗りする、削除する & 境界外転送前に個人データを「黒塗り」する。データ最小化の具体的実装手段。 \\
\W{remediate} & v. & 是正する、修復する & バイアス検出後は「是正」と継続監視がセット。検出だけで止まらない正解行動。 \\
\W{representative} & adj. & 代表的な、縮図となる & 評価セットは本番に「代表性」を持つこと。都合の良いサンプルは弱パターンの典型。 \\
\W{residency} & n. & （データの）居住性、所在 & EU 等、データ「所在」要件は性能最適化を上回る制約。GDPR 設計の正解鍵。 \\
\W{retention} & n. & 保持（期間） & データ「保持」を目的に必要な分だけに制限する。GDPR の核心で、最適化がこれに違反してはならない。 \\
\W{rubric} & n. & 評価基準表、採点ルーブリック & ブランド適合等の主観品質は LLM 審査の「ルーブリック」で定量化する。大規模評価の正解手法。 \\
\W{scrutiny} & n. & 精査、詳細な調査 & コードを書いた同一セッションでなく、独立レビューで「精査」を回復させる。AI コードレビューの正解肢。 \\
\W{sparse} & adj. & まばらな、希薄な & 10件程度の手動抽查は遅く「希薄」すぎて毎晩の問題を捉えられない。監視手法の不正解肢。 \\
\W{spectacle} & n. & 見世物、豪華さ & 最先端エージェント能力は「見世物」であり測定可能価値でない。初案件選定の不正解理由。 \\
\W{supersede} & v. & 取って代わる、廃止にする & 意思決定記録の覆しは元記録を「superseded（廃止）」扱いにして新記録を追加する（上書き・削除は不可）。ADR 運用の正解鍵。綴りは -s-（supercede は誤り）。 \\
\W{threshold} & n. & 閾値、基準値 & ローンチ許容「閾値」をあらかじめ合意しておく。「十分良い」は定義された数値「閾値」でのみ意味を持つ。 \\
\W{unambiguous} & adj. & 曖昧でない、明確な & 現行の決定を「明確」にするため superseded 状態を明示する。ADR 履歴の正しい扱い。 \\
\W{verbatim} & adv. & 逐語的に、一言一句そのまま & 規則を「逐語的に」繰り返すより構造化配置が有効。プロンプト改善の不正解肢。 \\
\W{warrant} & v. & 正当な根拠を与える、保証する & 能力拡大はより強い統制を「正当化する」。複数エージェント採用が「裏付けられる」かの判断語。 \\
\W{withhold} & v. & 差し控える、隠す & 決定者から「重要な情報を隠す」のは職業倫理違反。証拠開示の正解肢と対になる不正解理由。 \\

\end{longtable}

\end{document}
